\documentclass{article}
\usepackage[margin=2cm]{geometry}
\usepackage{amsmath}
\usepackage{graphicx}
\usepackage{booktabs}
\usepackage[utf8]{inputenc}
\usepackage{ragged2e}
\setlength{\emergencystretch}{3em}
\begin{document}
\sloppy

\noindent\hfill\includegraphics[height=1.8em]{assets/figure_header_logo.png}\par\noindent\hrule\vspace{4pt}

with zero corresponding to a NA value. Next we want to identify the
bivariate outliers by iterating the filter over all possible pairs of variables.
Consider a pair of variables $\mathbf { X ^ { [ j k ) } ~ = ~ \{ ( X . { i j. } X . { i k } ) \}, i ~ = ~ 1.... n}$ The ini-.
tial location and dispersion estimators are, respectively, the coordinate-wise
median and the.$2 \times 2$ sub-matrix $S ^ { ( j \pmb { k } ) }$ of the estimate computed by the$S$
generalized S-estimator on non-filtered data. Note that, this ensure the pos-.
itive definiteness property for and each$S$ $d \times d$ sub-matrix corresponding to
a subset of  variables. For bivariate points with no flagged components by.$^ { d }$
the univariate filter we compute the squared Mahalanobis distance.$\Delta . { i } ^ { ( \mathbf { j } \star ) }$ and
hence apply the bivariate filter, for all $1 < j < k < p$ At the end we want.
to identify the cells$( i, j )$ which have to be flagged as cell-wise outliers. The
procedure used for this purpose is described in ? and reported here. Let.


\begin{equation}
\begin{array} { c } { J = \{ ( i, j, k ) : \Delta _ { i } ^ { ( j k ) } ~ \mathrm { i s ~ f a g e d ~ a s ~ b i v a r i a t c ~ o u n l i e r } \} } \\ { \dots, \quad \dots, \quad \dots \dots \quad \dots \quad \dots \quad \dots \quad \dots \quad \dots \quad. } \end{array}
\end{equation}


be the set of triplets which identifies the pairs of cells flagged by the bivariate
filter in rows$i = 1, \dots, n$ For each cell in the data, we count the number$( i, j )$
of flagged pairs in the  th row in which the considered cell is involved:$i$


\begin{equation}
m _ { i j } = \# \{ k : ( i, j, k ) \in J \}.
\end{equation}


In absence of contamination,$n t . { i j }$ follows approximately a binomial distribution $B i r ( \sum . { k \neq j } \mathbf { U . { j k } }, \delta )$ whererepresents the overall proportion of cell-wise.$\diamondsuit$
outliers undetected by the univariate filter. Hence, we flag the cell $( i, j )$ if
$m . { i j } ~ > ~ c . { i j }$where$\epsilon . { i j }$is the 0.99-quantile of$B i n ( \sum . { k \neq j } \mathbf { U . { j k }, 0. 1 } )$.Finally, we
perform the -variate filter as described in subsection to the tull data matrix.$\boldsymbol { p }$
Detected observations (rows) are directly flagged as -variate (case-wise) out-$\rho$
liers. We denote the procedure based on univariate, bivariate and  -variate$\boldsymbol { p }$
filters, HS-UBPF.


\subsection*{0.1  A sequencing filtering procedure}


Suppose we would like to apply a sequence of filters with different dimension.$k ^ { * }$
$1 ~ \leq ~ d . { 1 } ~ \leq ~ d . { 2 } ~ \leq ~ \dots ~ \leq ~ d . { k } ~ \leq p$ For each $d . { i }$ $\dot { t } = 1, \dots, k$,the filter updates
the data matrix adding NA values to the -tuples identified as variate$t . { i }$ $a . { i }$
outliers. In this way, each filter applies only those  -tuples that have not$\boldsymbol { \omega } . { i }$
been flagged as outliers by the filters with lower dimension. Initial values for
each procedures rather than  would be obtained by applying the GSE to$\boldsymbol { u } . { 1 }$
the actual filtered values. This procedure aims to be a valid alternative to
that used in the presented HS-UBPF filter to perform a sequence of filters
with different dimensions. However, this is a preliminary idea, indeed it has
not been implemented yet.

\end{document}
